% BHU PYQ -- Manually transcribed & error-corrected
% Paper: BCH-213 : Specialised Accounts - I
% Exam: B.Com. (Hons.) Semester III Examination, 2016-17

\documentclass[11pt,a4paper]{article}
\usepackage[a4paper,top=2.5cm,bottom=2.5cm,left=2.8cm,right=2.8cm,headheight=14pt]{geometry}
\usepackage{amsmath,amssymb}
\usepackage[T1]{fontenc}
\usepackage[utf8]{inputenc}
\usepackage{lmodern,microtype,fancyhdr,enumitem,hyperref,lastpage,parskip}

\hypersetup{
    colorlinks=true,
    urlcolor=black,
    linkcolor=black,
    pdftitle={BCH-213: Specialised Accounts - I -- BHU B.Com. (Hons.) Sem III 2016-17}
}

\pagestyle{fancy}
\fancyhf{}
\renewcommand{\headrulewidth}{0.4pt}
\renewcommand{\footrulewidth}{0.4pt}
\fancyhead[L]{\small   \textit{Banaras Hindu University}}
\fancyhead[R]{\small   \textit{BCH-213: Specialised Accounts - I}}
\fancyfoot[C]{\small Page  \thepage\ of \pageref{LastPage}}


\newlist{parts}{enumerate}{2}
\setlist[parts,1]{label=(\roman*),leftmargin=*,itemsep=5pt,topsep=3pt}
\setlist[parts,2]{label=(\alph*),leftmargin=*,itemsep=3pt,topsep=2pt}

\newcommand{\pts}[1]{\hfill   \textit{[#1]}}


\begin{document}


\begin{center}
  {\large\bfseries BANARAS HINDU UNIVERSITY}\\[2pt]
  {\normalsize B.Com. (Hons.) --- Semester III Examination, 2016-17}\\[6pt]
  \rule{0.8\linewidth}{0.5pt}\\[6pt]
  {\large\bfseries Paper No. BCH-213: Specialised Accounts - I}\\[4pt]
  {\normalsize Time: 3 Hours\hfill Maximum Marks: 70}
\end{center}

\rule{\linewidth}{0.4pt}

\smallskip



\noindent   \textit{Instructions: Answer all questions. All questions carry equal marks (14 marks each).}

\bigskip




\noindent   \textbf{Question 1.}
What journal entries are passed in the books of Lessee and Lessor in connection with Royalty? Distinguish between royalty and rent.\pts{14}

\centerline{   \textbf{OR}}

On 01-01-2010 Mr. X acquired a coal mine on lease for 20 years at a royalty of Rs. 4 per ton of coal raised. Minimum rent is Rs. 48,000 in the first year and will increase by Rs. 2,000 every year till it reaches Rs. 60,000 per annum. Each year's shortworkings may be recoupable in the next two years only but on the condition that if full shortworkings could not be recouped in the next year, Mr. X will lose his right to recover 50\% of the unrecovered balance of shortworkings.

The output of the first six years is as follows:

\begin{center}
\small

\begin{tabular}{|c|c||c|c|}
\hline
   \textbf{Year} &   \textbf{Output (in tons)} &   \textbf{Year} &   \textbf{Output (in tons)} \\
\hline
2010 & 6,000  & 2013 & 40,000 \\
2011 & 16,000 & 2014 & 25,000 \\
2012 & 30,000 & 2015 & 35,000 \\
\hline
\end{tabular}
\end{center}

Prepare the necessary accounts in the books of Mr. X. Accounts are closed on 31st December every year.\pts{14}

\medskip\rule{\linewidth}{0.2pt}




\noindent   \textbf{Question 2.}
What is Hire Purchase System? Pass journal entries in case of instalment payment system in the books of the purchaser.\pts{14}

\centerline{   \textbf{OR}}

The cash price of a machine which is sold on hire purchase system on 1-1-2014 is Rs. 50,000. Rs. 10,000 are paid down and the balance in four equal quarterly instalments of Rs. 11,000 each. It should be noted that each of these instalments includes interest on outstanding balance. Find out the amount of interest included in each instalment and open the Hire Vendor's A/c in Hire Purchaser's books.\pts{14}

\pagebreak




\noindent   \textbf{Question 3.}
While preparing Departmental Trading and Profit \& Loss Account, why is the apportionment of indirect expenses necessary? Explain.\pts{14}

\centerline{   \textbf{OR}}

The following expenses were extracted from the books of Keshri for the year ended 31st December, 2015:


\begin{center}
\small

\begin{tabular}{|l|c|c|}
\hline
   \textbf{Particulars} &   \textbf{Department A (Rs.)} &   \textbf{Department B (Rs.)} \\
\hline
Opening Stock & 50,000 & 40,000 \\
Purchases & 6,50,000 & 9,10,000 \\
Sales & 10,00,000 & 15,00,000 \\
\hline
\end{tabular}
\end{center}

General expenses incurred for both the departments were Rs. 1,25,000. You are also supplied with the following information:

\begin{parts}
  \item Closing stock of Department A Rs. 1,00,000 including goods from Department B Rs. 20,000 at cost to Department A.
  \item Closing stock of Department B Rs. 2,00,000 including goods from Department A Rs. 30,000 at cost to Department B.
  \item Opening stock of Department A and Department B includes goods of the value of Rs. 10,000 and Rs. 15,000 taken from Department B and Department A respectively at cost to transferee department.
  \item The gross profit is uniform from year to year.
\end{parts}
You are required to prepare Departmental Trading \& P/L account for the year ended 31st December, 2015. Show your workings clearly.\pts{14}

\medskip\rule{\linewidth}{0.2pt}




\noindent   \textbf{Question 4.}
Why is branch accounting necessary? What are various kinds of branches? Discuss.\pts{14}

\centerline{   \textbf{OR}}

A firm has its head office at Delhi. Branches are at Agra and Varanasi. Accounting year is calendar year. On 31st December, 2015 Agra and Varanasi branch accounts in head office books show a debit balance of Rs. 30,000 and Rs. 45,000 respectively before taking the following into accounts:

\begin{parts}
  \item Rs. 3,000 was remitted from Varanasi to Agra branch.
  \item Goods valued at Rs. 2,000 were transferred from Agra to Varanasi under the instructions of head office.
  \item Agra branch collected Rs. 2,500 from an Agra customer of the head office.
  \item Varanasi branch paid Rs. 5,000 for certain goods purchased by the head office.
  \item Rs. 5,000 was remitted by the Agra branch to the head office on 30-12-15 but the head office received the same on 10th January, 2016.
  \item The Varanasi branch received on behalf of the head office Rs. 1,500 as dividend from a company in Varanasi.
  \item For the year 2015, Agra branch shows a net loss of Rs. 1,200 and Varanasi branch a net profit of Rs. 6,000.
\end{parts}
Prepare journal entries to record the above transactions in the head office books and write up the Agra and Varanasi branch accounts.\pts{14}

\pagebreak




\noindent   \textbf{Question 5.}
What do you understand by Human Resource Accounting? Discuss the Lev and Schwartz model of valuation of human resources.\pts{14}

\centerline{   \textbf{OR}}

Write notes on the following:\pts{14}

\begin{parts}
  \item Shortworkings
  \item Advantages of HR accounting
\end{parts}

\vfill
\rule{\linewidth}{0.4pt}

\begin{center}\small
\href{https://bhu-exam.vercel.app/}{Click here to visit our website}\\[4pt]
\href{https://me-aryan.vercel.app/}{Click here to visit our contributor}\\[4pt]
\href{https://quashan-me.vercel.app/}{Click here to visit our contributor}
\end{center}

\end{document}
